% answers.tex —— 1977 年·江苏扬州地区初试卷 · 参考答案与详细解析片段
% 由 tools/generate_projects.py 依来源总账自 1977/1977_数学真题_江苏扬州地区初试_试题.md + 1977/1977_数学真题_江苏扬州地区初试_解析正文.md + 1977/1977_数学真题_江苏扬州地区初试_答案要点.md 生成，勿手改（改源 MD 后重跑）


\begin{gktip}
固定源 PDF 仅含 2 页题面（原书页码 245—246），无解答可对照。以下各题均由编者独立推导，最终结论与答案要点逐题一致；原卷大题的中文数字编号「一、」至「十、」已转写为阿拉伯数字 \textbf{1．}至 \textbf{10．}。各题分值：第 1—4 题各 8 分，第 5—8 题各 10 分，第 9 题 16 分，第 10 题 12 分，合计 100 分。
\end{gktip}

\begin{answer}
\gksolhead{1}{}

\textbf{题目}（8 分）：

\Needspace{5\baselineskip}
计算：

\nopagebreak[4]
\[
8^{\frac{2}{3}}+\left(\frac{1}{2}\right)^{-2}-\log_2 1+\log_3 3-\left(\sqrt{2}-1\right)^{0}.
\]

\Needspace{4\baselineskip}
\textbf{解析}

逐项计算：$8^{\frac{2}{3}}=(2^{3})^{\frac{2}{3}}=2^{2}=4$；$\left(\dfrac{1}{2}\right)^{-2}=2^{2}=4$；$\log_2 1=0$；$\log_3 3=1$；$\sqrt{2}-1\ne0$，故 $\left(\sqrt{2}-1\right)^{0}=1$。

\Needspace{5\baselineskip}
于是

\nopagebreak[4]
\[
8^{\frac{2}{3}}+\left(\frac{1}{2}\right)^{-2}-\log_2 1+\log_3 3-\left(\sqrt{2}-1\right)^{0}
=4+4-0+1-1=8.
\]

\gkanswer{答案：$8$。}

\end{answer}
\begin{answer}
\gksolhead{2}{}

\textbf{题目}（8 分）：

求直线 $\sqrt{3}\,x+y-2=0$ 的斜率、倾斜角及在 $x$ 轴、$y$ 轴上的截距。

\Needspace{4\baselineskip}
\textbf{解析}

\Needspace{5\baselineskip}
将直线化成斜截式：

\nopagebreak[4]
\[
y=-\sqrt{3}\,x+2,
\]

故斜率 $k=-\sqrt{3}$。

\Needspace{5\baselineskip}
设倾斜角为 $\alpha$（$0^\circ\le\alpha<180^\circ$），则 $\tan\alpha=-\sqrt{3}$，故

\nopagebreak[4]
\[
\alpha=120^\circ.
\]

在 $x$ 轴上的截距：令 $y=0$，得 $x=\dfrac{2}{\sqrt{3}}=\dfrac{2\sqrt{3}}{3}$；在 $y$ 轴上的截距：令 $x=0$，得 $y=2$。

\gkanswer{答案：斜率 $k=-\sqrt{3}$，倾斜角为 $120^\circ$；在 $x$ 轴上的截距为 $\dfrac{2\sqrt{3}}{3}$，在 $y$ 轴上的截距为 $2$。}

\end{answer}
\begin{answer}
\gksolhead{3}{}

\textbf{题目}（8 分）：

把 $A$ 点的极坐标 $\left(4,\ \dfrac{3\pi}{4}\right)$ 化成直角坐标；把 $B$ 点的直角坐标 $(-\sqrt{3},\ 1)$ 化成极坐标。

\Needspace{4\baselineskip}
\textbf{解析}

\Needspace{5\baselineskip}
由极坐标化直角坐标公式 $x=\rho\cos\theta$，$y=\rho\sin\theta$：

\nopagebreak[4]
\[
x=4\cos\frac{3\pi}{4}=4\times\left(-\frac{\sqrt{2}}{2}\right)=-2\sqrt{2},
\qquad
y=4\sin\frac{3\pi}{4}=4\times\frac{\sqrt{2}}{2}=2\sqrt{2}.
\]

故 $A(-2\sqrt{2},\ 2\sqrt{2})$。

\Needspace{5\baselineskip}
对 $B(-\sqrt{3},\ 1)$：模长

\nopagebreak[4]
\[
\rho=\sqrt{(-\sqrt{3})^{2}+1^{2}}=\sqrt{4}=2.
\]

$\tan\theta=\dfrac{y}{x}=\dfrac{1}{-\sqrt{3}}=-\dfrac{\sqrt{3}}{3}$。因为 $x<0$、$y>0$，点 $B$ 在第二象限，故 $\theta=150^\circ=\dfrac{5\pi}{6}$。所以 $B$ 点的极坐标为 $\left(2,\ \dfrac{5\pi}{6}\right)$。

\gkanswer{答案：$A(-2\sqrt{2},\ 2\sqrt{2})$；$B$ 点的极坐标为 $\left(2,\ \dfrac{5\pi}{6}\right)$。}

\end{answer}
\begin{answer}
\gksolhead{4}{}

\textbf{题目}（8 分）：

不画图象求出函数 $y=2x^{2}-4x-1$ 的图象的顶点坐标、对称轴方程，并说明 $x$ 取何值时，函数取得极大值或极小值？这个极值是多少？

\Needspace{4\baselineskip}
\textbf{解析}

\Needspace{5\baselineskip}
配方：

\nopagebreak[4]
\[
y=2x^{2}-4x-1=2\left(x^{2}-2x\right)-1=2(x-1)^{2}-3.
\]

故图象（抛物线）的顶点坐标为 $(1,\ -3)$，对称轴方程为 $x=1$。

因二次项系数 $a=2>0$，抛物线开口向上，故当 $x=1$ 时函数取得极小值 $-3$，无极大值。

\gkanswer{答案：顶点 $(1,\ -3)$，对称轴 $x=1$；$x=1$ 时取得极小值 $-3$，无极大值。}

\end{answer}
\begin{answer}
\gksolhead{5}{}

\textbf{题目}（10 分）：

\Needspace{5\baselineskip}
已知 $x$ 是正的纯小数，化简

\nopagebreak[4]
\[
\sqrt{x^{2}-2+\frac{1}{x^{2}}}\cdot\left(x+\frac{1}{x}\right)\div\frac{2x^{-2}}{1-\dfrac{x^{2}}{1+x^{2}}}.
\]

\Needspace{4\baselineskip}
\textbf{解析}

\Needspace{5\baselineskip}
由 $0<x<1$ 得 $\dfrac{1}{x}>x>0$，故 $x-\dfrac{1}{x}<0$。于是

\nopagebreak[4]
\[
\sqrt{x^{2}-2+\frac{1}{x^{2}}}
=\sqrt{\left(x-\frac{1}{x}\right)^{2}}
=\left|x-\frac{1}{x}\right|
=\frac{1}{x}-x.
\]

\Needspace{5\baselineskip}
因此被除式为

\nopagebreak[4]
\[
\left(\frac{1}{x}-x\right)\left(x+\frac{1}{x}\right)
=\frac{1}{x^{2}}-x^{2}
=\frac{1-x^{4}}{x^{2}}.
\]

\Needspace{5\baselineskip}
除式的分母部分：

\nopagebreak[4]
\[
1-\frac{x^{2}}{1+x^{2}}=\frac{1}{1+x^{2}},
\]

\Needspace{5\baselineskip}
故整个除式为

\nopagebreak[4]
\[
\frac{2x^{-2}}{1-\dfrac{x^{2}}{1+x^{2}}}
=\frac{2}{x^{2}}\cdot(1+x^{2})
=\frac{2(1+x^{2})}{x^{2}}.
\]

\Needspace{5\baselineskip}
原式为

\nopagebreak[4]
\[
\frac{1-x^{4}}{x^{2}}\div\frac{2(1+x^{2})}{x^{2}}
=\frac{(1-x^{2})(1+x^{2})}{2(1+x^{2})}
=\frac{1-x^{2}}{2}.
\]

\gkanswer{答案：$\dfrac{1-x^{2}}{2}$。}

\end{answer}
\begin{answer}
\gksolhead{6}{}

\textbf{题目}（10 分）：

在 $\triangle ABC$ 中，$\angle A=60^\circ$，$\angle B=45^\circ$，$BC$ 边上的高 $AD=10$，求 $\triangle ABC$ 的面积（不查表，得数可保留根式）。

\Needspace{4\baselineskip}
\textbf{解析}

$\angle C=180^\circ-60^\circ-45^\circ=75^\circ$。设 $AD\perp BC$ 于 $D$。

在 $\mathrm{Rt}\triangle ABD$ 中，$\angle B=45^\circ$，故 $BD=AD=10$。

\Needspace{5\baselineskip}
在 $\mathrm{Rt}\triangle ACD$ 中，

\nopagebreak[4]
\[
DC=\frac{AD}{\tan C}=\frac{10}{\tan 75^\circ}.
\]

\Needspace{5\baselineskip}
由 $\tan 75^\circ=\tan(45^\circ+30^\circ)=2+\sqrt{3}$，得

\nopagebreak[4]
\[
DC=10(2-\sqrt{3}).
\]

\Needspace{5\baselineskip}
所以

\nopagebreak[4]
\[
BC=BD+DC=10+10(2-\sqrt{3})=10(3-\sqrt{3}).
\]

\Needspace{5\baselineskip}
于是

\nopagebreak[4]
\[
S_{\triangle ABC}=\frac{1}{2}\cdot BC\cdot AD
=\frac{1}{2}\times10(3-\sqrt{3})\times10
=50(3-\sqrt{3}).
\]

\gkanswer{答案：$50(3-\sqrt{3})$。}

\end{answer}
\begin{answer}
\gksolhead{7}{}

\textbf{题目}（10 分）：

\Needspace{5\baselineskip}
圆台上、下底面半径分别是 $2\ \mathrm{cm}$ 和 $6\ \mathrm{cm}$，它的侧面积是 $40\pi\ \mathrm{cm^{2}}$，求圆台体积。

\nopagebreak[4]
\[
\left[\;\left(V_{\text{\text{体}}}=\frac{1}{3}\pi h\left(R^{2}+r^{2}+Rr\right)\right)\;\right]
\]

\Needspace{4\baselineskip}
\textbf{解析}

\Needspace{5\baselineskip}
设母线长为 $l$。圆台侧面积

\nopagebreak[4]
\[
S_{\text{\text{侧}}}=\pi(R+r)l=\pi(6+2)l=40\pi,
\]

解得 $l=5$。

\Needspace{5\baselineskip}
高 $h$ 与母线、两底面半径差构成直角三角形：

\nopagebreak[4]
\[
h=\sqrt{l^{2}-(R-r)^{2}}=\sqrt{25-16}=3.
\]

\Needspace{5\baselineskip}
代入题给体积公式：

\nopagebreak[4]
\[
V_{\text{\text{体}}}=\frac{1}{3}\pi h\left(R^{2}+r^{2}+Rr\right)
=\frac{1}{3}\pi\times3\times(36+4+12)=52\pi\ (\mathrm{cm^{3}}).
\]

\gkanswer{答案：$52\pi\ \mathrm{cm^{3}}$。}

\end{answer}
\begin{answer}
\gksolhead{8}{}

\textbf{题目}（10 分）：

某工厂 1977 年一月份的产值是 10 万元，如果以后平均每月产值比前一个月增长 $10\%$，问到几月份就能完成今年总产值 135 万元的任务？（$\lg 1.1=0.0414$，$\lg 2.35=0.3711$）

\Needspace{4\baselineskip}
\textbf{解析}

\Needspace{5\baselineskip}
设从一月份算起到第 $n$ 个月累计产值达到 135 万元。第 $n$ 个月产值为 $10\times1.1^{\,n-1}$ 万元，前 $n$ 个月累计产值为等比数列和：

\nopagebreak[4]
\[
S_{n}=\frac{10\left(1.1^{\,n}-1\right)}{1.1-1}=100\left(1.1^{\,n}-1\right).
\]

\Needspace{5\baselineskip}
令 $S_{n}\ge135$，即

\nopagebreak[4]
\[
1.1^{\,n}\ge2.35.
\]

\Needspace{5\baselineskip}
两边取常用对数：

\nopagebreak[4]
\[
n\lg1.1\ge\lg2.35
\quad\Longrightarrow\quad
n\ge\frac{0.3711}{0.0414}\approx8.96.
\]

故 $n$ 最小取 $9$，即到九月份就能完成今年总产值 135 万元的任务。

\gkanswer{答案：到九月份（第 9 个月）就能完成全年总产值 135 万元的任务。}

\end{answer}
\begin{answer}
\gksolhead{9}{}

\textbf{题目}（16 分）：

已知一元二次方程 $2x^{2}+4(2e-1)x+4e^{2}-1=0$ 有两个相等的实数根，以 $e$ 作为椭圆的离心率，这个椭圆的中心在原点，焦点在 $x$ 轴上，半长轴的长是 10，求这个椭圆的方程。

\Needspace{4\baselineskip}
\textbf{解析}

\Needspace{5\baselineskip}
方程有两个相等的实数根，判别式为零：

\nopagebreak[4]
\[
\Delta=\left[4(2e-1)\right]^{2}-4\times2\times\left(4e^{2}-1\right)=0.
\]

\Needspace{5\baselineskip}
化简：

\nopagebreak[4]
\[
16(2e-1)^{2}-8(4e^{2}-1)=0
\quad\Longrightarrow\quad
2(2e-1)^{2}-(4e^{2}-1)=0,
\]

\nopagebreak[4]
\[
8e^{2}-8e+2-4e^{2}+1=0
\quad\Longrightarrow\quad
4e^{2}-8e+3=0.
\]

\Needspace{5\baselineskip}
解得

\nopagebreak[4]
\[
e=\frac{8\pm\sqrt{64-48}}{8}=\frac{8\pm4}{8},
\qquad
e=\frac{3}{2}\ \text{\text{或}}\ e=\frac{1}{2}.
\]

椭圆离心率须满足 $0<e<1$，故 $e=\dfrac{1}{2}$。

\Needspace{5\baselineskip}
由 $a=10$，$c=ae=5$，得

\nopagebreak[4]
\[
b^{2}=a^{2}-c^{2}=100-25=75.
\]

\Needspace{5\baselineskip}
所以椭圆的方程为

\nopagebreak[4]
\[
\frac{x^{2}}{100}+\frac{y^{2}}{75}=1.
\]

\gkanswer{答案：$\dfrac{x^{2}}{100}+\dfrac{y^{2}}{75}=1$。}

\end{answer}
\begin{answer}
\gksolhead{10}{}

\textbf{题目}（12 分）：

已知 $\triangle ABC$ 的三边 $a$、$b$、$c$ 成等差数列，$a>b>c$，$b$ 边的对角是 $\angle B$。求证：$\angle B$ 的度数小于 $60^\circ$。

\Needspace{4\baselineskip}
\textbf{解析}

\Needspace{5\baselineskip}
由 $a$、$b$、$c$ 成等差数列得

\nopagebreak[4]
\[
2b=a+c,
\qquad
b=\frac{a+c}{2}.
\]

\Needspace{5\baselineskip}
由余弦定理，

\nopagebreak[4]
\[
\cos B=\frac{a^{2}+c^{2}-b^{2}}{2ac}.
\]

\Needspace{5\baselineskip}
将 $b^{2}=\dfrac{(a+c)^{2}}{4}$ 代入：

\nopagebreak[4]
\[
a^{2}+c^{2}-b^{2}
=a^{2}+c^{2}-\frac{(a+c)^{2}}{4}
=\frac{3a^{2}-2ac+3c^{2}}{4},
\]

\Needspace{5\baselineskip}
故

\nopagebreak[4]
\[
\cos B=\frac{3a^{2}-2ac+3c^{2}}{8ac}.
\]

\Needspace{5\baselineskip}
因 $a>b>c>0$，故 $a\ne c$。由 $3a^{2}+3c^{2}>6ac$（严格不等式），得

\nopagebreak[4]
\[
3a^{2}-2ac+3c^{2}>4ac
\quad\Longrightarrow\quad
\cos B>\frac{4ac}{8ac}=\frac{1}{2}.
\]

\Needspace{5\baselineskip}
又 $B$ 为三角形内角，$0^\circ<B<180^\circ$，$\cos B$ 在此区间上单调递减，故

\nopagebreak[4]
\[
\angle B<60^\circ.
\]

\gkanswer{答案：见解析版}

\end{answer}